\section{Results}
\label{sec:results}

Unless stated otherwise, bounds are for stage~D (width, thickness, strip angle, etch depth, rib angle, $\nsi$ and
$\nox$ unknown) with realistic priors, $\sigma_\lambda=\qty{1}{pm}$, at derivative step $h$, and the bracket gives
the range over $h/2$, $h$ and $2h$. Table~\ref{tab:summary} collects the main sets.

\begin{table}[!t]
\caption{Strip sidewall-angle CRLB $\sth$ for the main test-structure sets (stage~D, realistic priors). Ranges are
over the derivative steps $h/2$, $h$, $2h$; ``MC'' is the standard deviation of 2000 nonlinear fits at
\qty{1}{pm}.}
\label{tab:summary}
\centering
\setlength{\tabcolsep}{3pt}
\begin{tabular}{@{}lcccc@{}}
\toprule
Set & \multicolumn{2}{c}{$\sth$ at \qty{1}{pm} (\si{\degree})} & \qty{0.1}{pm} & MC \\
 & $h$ & range & (\si{\degree}) & (\si{\degree}) \\
\midrule
Strips only, any number & $\infty$ & -- & $\infty$ & -- \\
Strips + one rib (\qty{70}{nm} slab) & 25.19 & -- & -- & -- \\
Strips + rib 300, 500 (\qty{50}{nm}) & 0.3891 & 0.3891--0.3997 & 0.04077 & 0.3976 \\
Best set of 8 (\qty{90}{nm}) & 0.139 & 0.1007--0.1473 & 0.0139 & 0.1323 \\
\bottomrule
\end{tabular}
\end{table}

\subsection{Strips are blind to the angle}

Fig.~\ref{fig:sections} shows the mechanism. For the strip the wall weight $W(y)$ is symmetric about mid-height and
its centroid sits at $y_c=t/2$; for the rib the weight is concentrated near the slab corner and the centroid lies below
the middle of the wall. The femwell derivatives follow~\eqref{eq:R}. Over all ten strips, in TE and TM and for both
$\neff$ and $\ngr$, the ratio $(\partial n/\partial\theta)/(\partial n/\partial w)$ lies between \qty{3.8396}{nm} and
\qty{3.8725}{nm} per degree, against $t\tan(\qty{1}{\degree})=\qty{3.8401}{nm}$; the largest deviation is
\qty{0.84}{\percent}. The closed-form $R$ computed from the mode-field centroid matches the femwell ratio within
\qty{0.79}{\percent} for the strips, a median of \qty{1.25}{\percent} for the ribs, and at worst \qty{5.4}{\percent}.
For ribs, $R$ lies between 0.831 and 0.873 on the \qty{50}{nm} slab and between 0.750 and 0.799 on the
\qty{90}{nm} slab, and it falls as the slab thickens (Fig.~\ref{fig:ratio}). With the symmetry imposed, the best set of
strips of any size from one to eight has an infinite bound: strips alone cannot determine the angle.

\begin{figure*}[!t]
\centering
\includegraphics{fig1_sections}
\caption{Why a strip cannot see its sidewall angle. (a)~Oxide-clad strip, \qty{500}{nm}~$\times$~\qty{220}{nm}, and
(b)~rib, \qty{500}{nm} wide on a \qty{90}{nm} slab, both TE at $\lambda=\qty{1.55}{\micro m}$, with the transverse
intensity $|\mathbf{E}_t|^2$, the nominal wall (solid) and a tilted wall (dashed; \qty{8}{\degree}, exaggerated for
visibility) pivoting about the bottom corner (dot). The femwell angle/width ratio $R$ is printed in each panel.
(c)~First-order wall-perturbation weight $W$ along the right wall versus height above the pivot, normalized to its
maximum, with the centroids $y_c$ (dashed). The strip weight is symmetric about mid-height, so a tilt is exactly a width
change ($R=1$); the rib weight sits low on the wall, so $R<1$.}
\label{fig:sections}
\end{figure*}

\begin{figure}[!t]
\centering
\includegraphics{fig2_ratio}
\caption{Angle/width ratio $R$ of the TE ribs versus slab thickness for four rib widths (femwell derivatives at step
$h$; $\lambda=\qty{1.55}{\micro m}$, $t=\qty{220}{nm}$, oxide-clad). Open symbols: closed form
$R=(y_c-y_\mathrm{pivot})/(h/2)$ from the mode field at the \qty{50}{nm} and \qty{90}{nm} slabs. The line at $R=1$ is
the strip value, set by symmetry. Every rib breaks the strip degeneracy, and thinner slabs break it least.}
\label{fig:ratio}
\end{figure}

\subsection{One rib against unknown etch parameters}

Fig.~\ref{fig:stages} follows a rib on a \qty{70}{nm} slab through the stages. If the rib shares the strip angle
(stage~B) and only its depth is unknown, one rib width already gives \qty{0.357}{\degree} and two widths give
\qty{0.00365}{\degree}. Giving the rib its own angle (stage~C) changes the picture. The single rib then leaves the strip
angle unidentifiable, as the count in Section~\ref{sec:structures} predicts, while the two-width set loses only part of
its precision and gives \qty{0.448}{\degree} without priors. Adding $\nsi$, $\nox$ (stage~D) and temperature (stage~E)
moves the two-width bound to \qty{0.481}{\degree} without priors and to \qty{0.452}{\degree} and \qty{0.455}{\degree}
with realistic priors.

\begin{figure}[!t]
\centering
\includegraphics{fig3_stages}
\caption{Strip sidewall-angle CRLB as process unknowns are added (B: etch depth, rib angle tied to the strip angle;
C: separate rib angle; D: silicon and oxide indices; E: temperature), for strips alone (crosses), strips plus one
\qty{500}{nm} rib (diamonds) and strips plus ribs of \qty{500}{nm} and \qty{300}{nm} (circles), all on a \qty{70}{nm}
slab, without priors (open) and with realistic priors (filled), at $\sigma_\lambda=\qty{1}{pm}$. Cases with a singular
Fisher matrix are drawn in the band marked $\infty$. A second width of the same rib is what keeps the angle
identifiable once the partial etch has its own angle.}
\label{fig:stages}
\end{figure}

\subsection{Priors do not rescue a single rib}

With a single rib whose etch angle is independent of the strip's (stage~C, \qty{70}{nm} slab, \qty{1}{pm}), the strip
angle is unidentifiable without priors. Prior knowledge improves this only partially: realistic priors give
$\sth=\qty{25.18}{\degree}$, knowing the etch depth to \qty{0.01}{nm} (far beyond what a separate thickness measurement
provides) gives \qty{5.03}{\degree}, and tying the rib and strip angles to within \qty{1}{\degree} gives \qty{1.25}{\degree}. Two widths of the same rib reach \qty{0.448}{\degree} at the same stage without any of these
assumptions. In stage~D with realistic priors, the single rib gives between \qty{24.4}{\degree} and \qty{45.4}{\degree}
depending on the slab; these values are set by the priors, not by the data.

\subsection{Multi-width rib sets}

The two-width set, strip TE rings at the five widths plus ribs of \qty{300}{nm} and \qty{500}{nm} on a \qty{50}{nm} slab,
gives $\sth=\qty{0.389}{\degree}$ (\qtyrange{0.389}{0.400}{\degree} over steps). It is also insensitive to
derivative errors: on a fixed mesh the bound is \qty{0.394}{\degree}, a \qty{1.3}{\percent} change, and random errors of
$\pm\qty{0.3}{\percent}$ on every \qty{300}{nm}-rib derivative spread it only between \qty{0.364}{\degree} and \qty{0.420}{\degree} (5th to 95th percentile).

An exhaustive search over all sets of up to eight structures that use one partial etch finds a lower bound with more
rib widths and some TM strips. The best set of eight, TE strips of \qty{400}{nm} and \qty{600}{nm}, TM strips of
\qty{400}{nm} and \qty{600}{nm} and ribs of 300, 400, 500 and \qty{600}{nm} on a \qty{90}{nm} slab, gives
$\sth=\qty{0.139}{\degree}$ (\qtyrange{0.101}{0.147}{\degree} over steps). The best sets of four to seven structures
give 0.435, 0.308, 0.153 and \qty{0.145}{\degree}. This set is more sensitive to derivative accuracy than the
two-width set. With fixed-mesh derivatives for the \qty{300}{nm} rib its bound becomes \qty{0.118}{\degree} (\qty{-15.3}{\percent}), the fixed-mesh and remeshed derivatives differing by at most \qty{0.78}{\percent}, and random
$\pm\qty{0.3}{\percent}$ errors spread it between \qty{0.109}{\degree} and \qty{0.189}{\degree}. We therefore quote
\qtyrange{0.12}{0.14}{\degree} for this set, with a derivative-accuracy uncertainty of about \qty{15}{\percent}. Without the
index unknowns (stage~C, no priors) the best set of eight gives \qty{0.124}{\degree}.

A greedy search that adds one structure at a time reached the exhaustive optimum in 25 of the 32 stage, prior and slab
cases we compared and fell short in 7, with a bound up to a factor of 1.12 larger. At the best-case \qty{0.1}{pm} wavelength precision
both bounds shrink tenfold, to \qty{0.0408}{\degree} for the two-width set and \qty{0.0139}{\degree} for the best set;
they scale linearly with $\sigma_\lambda$ (\qty{0.196}{\degree} and \qty{0.775}{\degree} for the two-width set at 0.5
and \qty{2}{pm}).

\subsection{Monte Carlo fits reach the bound}

\begin{figure*}[!t]
\centering
\includegraphics{fig4_montecarlo}
\caption{Monte Carlo fits of 2000 simulated noisy measurement sets per set at $\sigma_\lambda=\qty{1}{pm}$
(stage~D, realistic priors, bounded nonlinear least squares). (a)~Histograms of the fitted strip angle for strip TE
rings at five widths plus one \qty{500}{nm} rib (grey) and plus ribs of \qty{300}{nm} and \qty{500}{nm} (orange), both on
a \qty{50}{nm} slab; dotted lines mark the $\pm\qty{3}{\degree}$ fit box. With one rib the fits spread over the whole box
along the width/angle valley. (b)~Fitted-angle histograms for the two-width set and the best set of eight
(\qty{90}{nm} slab) with Gaussians of the CRLB width (lines): with two or more rib widths the fit scatter matches the
bound.}
\label{fig:mc}
\end{figure*}

Fig.~\ref{fig:mc} compares the nonlinear fits with the bound. For the two-width set the fitted angles have a standard
deviation of \qty{0.398}{\degree} (\qtyrange{0.385}{0.410}{\degree}, 95\% bootstrap interval), 1.02 times the CRLB,
with a bias of \qty{-0.015 \pm 0.0089}{\degree}; none of the 2000 fits failed, ended ambiguous or reached the box edge.
For the best set the spread is \qty{0.132}{\degree} (\qtyrange{0.125}{0.140}{\degree}), 0.952 times the CRLB, with a
bias of \qty{0.0075 \pm 0.003}{\degree}; 4 fits ended at the box edge and none failed. The spread is slightly below the
femwell CRLB because the Monte Carlo uses the surrogate model, whose own bound at nominal is \qty{0.1249}{\degree}: its
\qty{300}{nm}-rib slopes are a smooth fit through the remeshing scatter. Measured against that bound the fits are
slightly less efficient than the CRLB, as expected for a nonlinear estimator. With one rib the fits find the valley predicted by the
symmetry argument. The linearized bound is \qty{29.52}{\degree}, so the fit box of $\pm\qty{3}{\degree}$ confines the
estimates rather than the data: 1190 of the 2000 fits end at the edge of the box, 12 are ambiguous and 20 fail, and the
fitted widths and angles lie on a line of slope \qty{3.842}{nm} per degree with a correlation of 0.999998, against
$t\tan(\qty{1}{\degree})=\qty{3.8401}{nm}$ (Fig.~\ref{fig:mc}(a)).

\subsection{Slab thickness}

For the two-width set the bound depends strongly on the slab (Fig.~\ref{fig:slab}). It is lowest at \qty{50}{nm}
(\qty{0.389}{\degree}) and rises slowly toward thinner slabs, to \qty{0.403}{\degree} at \qty{40}{nm} and
\qty{0.414}{\degree} at \qty{30}{nm}, and quickly toward thicker ones, to \qty{0.597}{\degree} at \qty{90}{nm} and
\qty{4.089}{\degree} at \qty{150}{nm}. Among the slabs simulated the optimum is therefore \qty{50}{nm}, not
the thinnest slab. Thinner slabs bring $R$ closer to 1 (Fig.~\ref{fig:ratio}), but we have not separated that effect
from the other ways the slab changes the two rib modes.

\begin{figure}[!t]
\centering
\includegraphics{fig5_slab}
\caption{CRLB of the two-width set (strip TE rings at 400--\qty{600}{nm} plus ribs of \qty{300}{nm} and
\qty{500}{nm}) versus slab thickness, stage~D, realistic priors, $\sigma_\lambda=\qty{1}{pm}$, $\lambda=\qty{1.55}{\micro m}$.
Bars span the derivative steps $h/2$, $h$, $2h$. The bound is lowest at \qty{50}{nm} and grows tenfold by
\qty{150}{nm}.}
\label{fig:slab}
\end{figure}

\subsection{Cost of ignoring the angle}

When strip rings are fitted with the angle held at its nominal value, a real \qty{1}{\degree} tilt is absorbed as a width
change of \qty{-3.84}{nm}. Using that width to predict a \qty{500}{nm} rib ring on a \qty{90}{nm} slab, whose angle
sensitivity differs ($R<1$), mispredicts its resonance by $\delta\lambda=\lambda\,\delta\neff/\ngr=\qty{-711}{pm}$ per
degree of unmodeled tilt.
